% year: תשע"ט
% type: מבחן
% semester: א
% moed: מיוחד
% number: 4
% points: 20
% categories: func-analysis, derivatives
% source: exams/מבחנים/תשעט/מועד מיוחד תשעט.pdf

\begin{question}
חקרו את הפונקציה $f(x)=\dfrac{e^{2x-1}}{x^2}$ וסרטטו את הגרף שלה.
צריך למצוא: תחום הגדרה, אסימפטוטות, תחומי עליה וירידה, נקודות קיצון, נקודות חיתוך של הגרף עם הצירים, תחומי קמירות וקעירות.
\end{question}

\begin{hint}
$f'(x)=\frac{2e^{2x-1}(x-1)}{x^3}$. שימו לב שהסימן של $x^3$ משתנה ב-$x=0$.
\end{hint}

\begin{solution}
\textbf{תחום הגדרה:} $x\neq 0$, כלומר $(-\infty,0)\cup(0,\infty)$.

\textbf{חיתוך עם הצירים:} $x=0$ אינו בתחום, ולכן אין חיתוך עם ציר $y$. מכיוון ש-$e^{2x-1}>0$ ו-$x^2>0$, מתקיים $f(x)>0$ לכל $x$ בתחום, ולכן אין חיתוך עם ציר $x$.

\textbf{אסימפטוטות:}
\begin{itemize}
  \item \emph{אנכית:} $\lim_{x\to0^\pm}\frac{e^{2x-1}}{x^2} = \frac{e^{-1}}{0^+} = +\infty$, לכן $x=0$ אסימפטוטה אנכית (הגרף עולה ל-$+\infty$ משני הצדדים).
  \item \emph{ב-$-\infty$:} $e^{2x-1}\to0$ ו-$\frac1{x^2}\to0$, לכן $\lim_{x\to-\infty}f(x)=0$, ו-$y=0$ אסימפטוטה אופקית משמאל.
  \item \emph{ב-$+\infty$:} $\frac{f(x)}{x}=\frac{e^{2x-1}}{x^3}\to\infty$ (אקספוננט גובר על כל חזקה, למשל לפי לופיטל שלוש פעמים), ולכן אין אסימפטוטה אופקית או משופעת בצד ימין, ו-$f(x)\to+\infty$.
\end{itemize}

\textbf{נגזרת ראשונה:} לפי כלל המנה,
\[ f'(x) = \frac{2e^{2x-1}x^2 - e^{2x-1}\cdot 2x}{x^4} = \frac{2e^{2x-1}(x-1)}{x^3}. \]
$f'(x)=0 \iff x=1$. הסימן של $f'$ נקבע ע"י $\frac{x-1}{x^3}$:
\begin{itemize}
  \item $x<0$: $x-1<0$, $x^3<0$ ולכן $f'>0$ — $f$ \textbf{עולה} ב-$(-\infty,0)$.
  \item $0<x<1$: $x-1<0$, $x^3>0$ ולכן $f'<0$ — $f$ \textbf{יורדת} ב-$(0,1)$.
  \item $x>1$: $f'>0$ — $f$ \textbf{עולה} ב-$(1,\infty)$.
\end{itemize}
לכן ב-$x=1$ יש \textbf{מינימום מקומי}: $f(1)=\frac{e^{1}}{1}=e$, כלומר הנקודה $(1,e)$. אין נקודות קיצון נוספות ($x=0$ אינו בתחום).

\textbf{נגזרת שנייה:} נגזור את $f'(x)=2e^{2x-1}\cdot\frac{x-1}{x^3}$:
\[ \left(\frac{x-1}{x^3}\right)' = \frac{x^3-3x^2(x-1)}{x^6} = \frac{-2x+3}{x^4}, \]
\[ f''(x) = 2e^{2x-1}\left(\frac{2(x-1)}{x^3}+\frac{3-2x}{x^4}\right) = \frac{2e^{2x-1}\left(2x^2-4x+3\right)}{x^4}. \]
לטרינום $2x^2-4x+3$ דיסקרימיננטה $16-24<0$ ומקדם מוביל חיובי, ולכן הוא חיובי לכל $x$. גם $e^{2x-1}>0$ ו-$x^4>0$. לכן $f''(x)>0$ לכל $x\ne0$:
$f$ \textbf{קמורה} (קמורה כלפי מעלה, $\cup$) בכל אחד מהקטעים $(-\infty,0)$ ו-$(0,\infty)$, ואין נקודות פיתול.

\textbf{סרטוט:} משמאל הגרף מתחיל קרוב לציר $x$ (מעליו), עולה וקמור ושואף ל-$+\infty$ כאשר $x\to0^-$. מימין ל-$0$ הגרף יורד מ-$+\infty$ עד המינימום $(1,e)$ ואז עולה במהירות אקספוננציאלית ל-$+\infty$. כל הגרף מעל ציר $x$.
\end{solution}
